\sectionheading{5.16 প্রোগ্রাম ডিজাইন মডেল (Program Design Model)}

কোনো সমস্যা সমাধানের জন্য শুধু Program লিখলেই হয় না;
Program তৈরির কাজটি \textbf{কোন পদ্ধতিতে পরিকল্পনা, উন্নয়ন,
পরীক্ষা ও বাস্তবায়ন করা হবে}, সেটিও গুরুত্বপূর্ণ।

Program তৈরির এই সুশৃঙ্খল পদ্ধতি বা কাঠামোকেই সাধারণভাবে
\textbf{Program Design Model} বলা হয়।

%================================================
% Definition
%================================================
\booksubheading{Program Design Model কাকে বলে?}

\begin{definitionbox}

\textbf{\color{darkTeal}সংজ্ঞা}

Program তৈরির বিভিন্ন ধাপকে একটি নির্দিষ্ট পদ্ধতি ও কাঠামো
অনুসারে সাজিয়ে Program উন্নয়নের যে Model তৈরি করা হয়,
তাকে \textbf{Program Design Model} বলে।

\end{definitionbox}

\textbf{সহজভাবে মনে রাখবে—}

\begin{center}

\textbf{
Program কীভাবে ধাপে ধাপে তৈরি হবে
\(\;=\;\)
Program Design Model
}

\end{center}

%================================================
% Importance
%================================================
\booksubheading{Program Design Model কেন প্রয়োজন?}

একটি ভালো Design Model ব্যবহার করলে—

\begin{compacttwocolpoints}

\item Program তৈরির কাজ পরিকল্পিতভাবে করা যায়।

\item বড় Problem-কে সহজে পরিচালনা করা যায়।

\item ভুল ও জটিলতা কমানো যায়।

\item Testing ও Modification সহজ হয়।

\item Program Development-এর অগ্রগতি সহজে বোঝা যায়।

\item Team-এর একাধিক সদস্যের কাজ সমন্বয় করা সুবিধাজনক হয়।

\end{compacttwocolpoints}

%================================================
% Process Model Intro
%================================================
\booksubheading{Program Development-এর Process Model}

Program বা Software তৈরির ধাপগুলো কীভাবে সম্পন্ন হবে তার ভিত্তিতে
বিভিন্ন ধরনের Process Model ব্যবহৃত হয়।

HSC পর্যায়ে প্রধান ধারণাগুলো হলো—

%================================================
% Linear
%================================================
\booksubheading{১. রৈখিক প্রক্রিয়া মডেল (Linear Process Model)}

\begin{definitionbox}

যে Model-এ Program Development-এর একটি ধাপ শেষ হওয়ার পর
ধারাবাহিকভাবে পরবর্তী ধাপে যাওয়া হয়, তাকে
\textbf{Linear Process Model} বলা হয়।

\end{definitionbox}

\textbf{সহজ ধারণা—}

\begin{center}
\begin{adjustbox}{max width=0.98\linewidth}
\fcolorbox{darkBlue}{softBlue}{
    \parbox[c][8mm][c]{0.14\linewidth}{
        \centering\textbf{ধাপ--১}
    }
}
\(\longrightarrow\)
\fcolorbox{mainTeal}{softGreen}{
    \parbox[c][8mm][c]{0.14\linewidth}{
        \centering\textbf{ধাপ--২}
    }
}
\(\longrightarrow\)
\fcolorbox{darkBlue}{softBlue}{
    \parbox[c][8mm][c]{0.14\linewidth}{
        \centering\textbf{ধাপ--৩}
    }
}
\(\longrightarrow\)
\fcolorbox{mainTeal}{softGreen}{
    \parbox[c][8mm][c]{0.14\linewidth}{
        \centering\textbf{ধাপ--৪}
    }
}
\end{adjustbox}
\end{center}

অর্থাৎ কাজ সাধারণত একটি নির্দিষ্ট ক্রমে এগোয়।

\textbf{উদাহরণ:}
Waterfall Model এবং V-Model।

%================================================
% Waterfall
%================================================
\booksubheading{Waterfall Model}

Waterfall Model-এ Development-এর ধাপগুলো সাধারণত ওপর থেকে নিচে
ধারাবাহিকভাবে সম্পন্ন হয়।

\begin{center}

\textbf{
Requirement Analysis
\(\rightarrow\)
Design
\(\rightarrow\)
Implementation
\(\rightarrow\)
Testing
\(\rightarrow\)
Deployment
\(\rightarrow\)
Maintenance
}

\end{center}

এর মূল ধারণা হলো—
\textbf{একটি ধাপ শেষ করে পরবর্তী ধাপে যাওয়া}।

%================================================
% Iterative
%================================================
\booksubheading{২. পুনরাবৃত্তিমূলক প্রক্রিয়া মডেল (Iterative Process Model)}

\begin{definitionbox}

যে Model-এ Program বা System একবারে সম্পূর্ণ তৈরি না করে
\textbf{বারবার Development Cycle-এর মধ্য দিয়ে উন্নত করা হয়},
তাকে \textbf{Iterative Process Model} বলে।

\end{definitionbox}

\begin{center}

\textbf{
Design
\(\rightarrow\)
Develop
\(\rightarrow\)
Test
\(\rightarrow\)
Improve
\(\rightarrow\)
আবার Development
}

\end{center}

প্রথম Version সম্পূর্ণ নিখুঁত না হলেও পরবর্তী Iteration-এ
নতুন Feature যোগ বা প্রয়োজনীয় পরিবর্তন করা যায়।

\textbf{উদাহরণ:} Spiral Model।

%================================================
% Parallel
%================================================
\booksubheading{৩. সমান্তরাল প্রক্রিয়া মডেল (Parallel Process Model)}

\begin{definitionbox}

যে Model-এ Development-এর বিভিন্ন কাজ বা অংশ প্রয়োজন অনুযায়ী
\textbf{একই সময় বা পরস্পরের সঙ্গে সমন্বয় করে} সম্পন্ন করা যায়,
তাকে \textbf{Parallel Process Model} বলা হয়।

\end{definitionbox}

এক্ষেত্রে একটি কাজ পুরোপুরি শেষ হওয়ার জন্য সবসময় অন্য কাজকে
অপেক্ষা করতে হয় না।

\begin{center}

\textbf{
একাধিক Development Task
\(\longrightarrow\)
সমান্তরালে কাজ
\(\longrightarrow\)
সমন্বয়
\(\longrightarrow\)
সম্পূর্ণ Program
}

\end{center}

\textbf{উদাহরণ:} Unified Process Model।

%================================================
% Process Model Comparison
%================================================
\booksubheading{Process Model-এর তুলনা}

\begin{tabularx}{\linewidth}{|>{\bfseries}p{0.19\linewidth}|Y|Y|Y|}
\hline

\rowcolor{softBlue}
\multicolumn{1}{|c|}{\textbf{বিষয়}}
&
\multicolumn{1}{c|}{\textbf{Linear}}
&
\multicolumn{1}{c|}{\textbf{Iterative}}
&
\multicolumn{1}{c|}{\textbf{Parallel}}
\\
\hline

কাজের ধরন
&
ধারাবাহিক
&
পুনরাবৃত্তিমূলক
&
সমান্তরাল
\\
\hline

ধাপ
&
একটির পর একটি
&
Development Cycle বারবার হতে পারে
&
কিছু কাজ একই সময়ে চলতে পারে
\\
\hline

পরিবর্তন
&
পরে পরিবর্তন তুলনামূলক কঠিন
&
পরিবর্তন ও উন্নয়ন তুলনামূলক সহজ
&
বিভিন্ন অংশ একসঙ্গে উন্নয়ন সম্ভব
\\
\hline

উদাহরণ
&
Waterfall, V-Model
&
Spiral Model
&
Unified Process Model
\\
\hline

\end{tabularx}

%================================================
% Programming Model
%================================================
\booksubheading{Programming Model-এর ধারণা}

Program-এর Code কীভাবে সংগঠিত ও তৈরি করা হবে তার ভিত্তিতেও
বিভিন্ন Programming Model রয়েছে।

এ অধ্যায়ের জন্য গুরুত্বপূর্ণ Modelগুলো হলো—

%================================================
% Structured
%================================================
\booksubheading{১. Structured Programming}

Program-কে সুশৃঙ্খলভাবে বিভিন্ন অংশে ভাগ করে
\textbf{Sequence, Selection এবং Repetition}-এর মতো
Control Structure ব্যবহার করে Program তৈরি করার পদ্ধতিকে
\textbf{Structured Programming} বলা হয়।

\textbf{মূল ধারণা:}
Program-কে ছোট, সুশৃঙ্খল ও সহজ অংশে সাজানো।

%================================================
% OOP
%================================================
\booksubheading{২. Object Oriented Programming (OOP)}

যে Programming পদ্ধতিতে Data এবং তার সঙ্গে সম্পর্কিত Operation-কে
\textbf{Object}-এর ভিত্তিতে সংগঠিত করে Program তৈরি করা হয়,
তাকে \textbf{Object Oriented Programming} বলে।

OOP-এর কয়েকটি গুরুত্বপূর্ণ ধারণা—

\begin{compacttwocolpoints}

\item Object
\item Class
\item Encapsulation
\item Inheritance
\item Polymorphism

\end{compacttwocolpoints}

এখানে এগুলোর নাম ও মূল ধারণা মনে রাখলেই যথেষ্ট।

%================================================
% Visual
%================================================
\booksubheading{৩. Visual Programming}

যে Programming পদ্ধতিতে Graphical Interface, Form, Button,
Menu ইত্যাদি Visual Element ব্যবহার করে Program তৈরি করা যায়,
তাকে \textbf{Visual Programming} বলে।

এতে Text Code-এর পাশাপাশি Graphical Element ব্যবহার করেও
Program তৈরির বিভিন্ন কাজ করা যায়।

%================================================
% Event Driven
%================================================
\booksubheading{৪. Event Driven Programming}

যে Programming পদ্ধতিতে কোনো \textbf{Event} সংঘটিত হলে
সেই Event অনুযায়ী Program-এর নির্দিষ্ট Code Execute হয়,
তাকে \textbf{Event Driven Programming} বলে।

Event-এর উদাহরণ—

\begin{itemize}

\item Mouse Click
\item Key Press
\item Button Click

\end{itemize}

\begin{examplebox}{সহজ উদাহরণ}

একটি Application-এ কোনো Button-এ Click করলে Result প্রদর্শিত হলো।

এখানে \textbf{Button Click} হলো Event এবং সেই Event-এর কারণে
নির্দিষ্ট Code Execute হলো।

\end{examplebox}

%================================================
% Programming Model Comparison
%================================================
\booksubheading{Programming Model-এর তুলনা}

\begin{tabularx}{\linewidth}{|>{\bfseries}p{0.33\linewidth}|Y|}
\hline

\rowcolor{softBlue}
\multicolumn{1}{|c|}{\textbf{Model}}
&
\multicolumn{1}{c|}{\textbf{মূল ধারণা}}
\\
\hline

Structured Programming
&
Program-কে সুশৃঙ্খল Control Structure-এ সাজানো।
\\
\hline

Object Oriented Programming
&
Object ও Class-এর ভিত্তিতে Program তৈরি।
\\
\hline

Visual Programming
&
Graphical Element ব্যবহার করে Program তৈরি।
\\
\hline

Event Driven Programming
&
Event ঘটলে সংশ্লিষ্ট Code Execute করা।
\\
\hline

\end{tabularx}

%================================================
% Process vs Programming Model
%================================================
\booksubheading{Process Model ও Programming Model-এর পার্থক্য}

দুটি বিষয় দেখতে কাছাকাছি হলেও এক নয়।

\begin{tabularx}{\linewidth}{|>{\bfseries}p{0.25\linewidth}|Y|Y|}
\hline

\rowcolor{softBlue}
\multicolumn{1}{|c|}{\textbf{বিষয়}}
&
\multicolumn{1}{c|}{\textbf{Process Model}}
&
\multicolumn{1}{c|}{\textbf{Programming Model}}
\\
\hline

মূল প্রশ্ন
&
Program তৈরির কাজ কীভাবে এগোবে?
&
Program-এর Code কীভাবে সংগঠিত হবে?
\\
\hline

উদাহরণ
&
Linear, Iterative, Parallel
&
Structured, OOP, Visual, Event Driven
\\
\hline

গুরুত্ব
&
Development Process
&
Programming Style/Structure
\\
\hline

\end{tabularx}

%================================================
% Final Reminder
%================================================
\begin{rememberbox}

\textbf{\color{deepOrange}মনে রাখবে}

\begin{itemize}

\item \textbf{Program Design Model} Program Development-কে
পরিকল্পিত ও সুশৃঙ্খল করে।

\item \textbf{Linear Model} →
ধাপগুলো ধারাবাহিকভাবে সম্পন্ন হয়।

\item \textbf{Waterfall ও V-Model} →
Linear Model-এর উদাহরণ।

\item \textbf{Iterative Model} →
Development Cycle বারবার সম্পন্ন করে Program উন্নত করা হয়।

\item \textbf{Spiral Model} →
Iterative Model-এর উদাহরণ।

\item \textbf{Parallel Model} →
একাধিক Development কাজ একই সময়ে চলতে পারে।

\item \textbf{Unified Process Model} →
Parallel Process Model-এর উদাহরণ।

\item Programming Model-এর গুরুত্বপূর্ণ ধরন—
\textbf{Structured, Object Oriented, Visual ও Event Driven Programming}।

\item দ্রুত পার্থক্য—

\begin{center}

\textbf{
Process Model
=
Program তৈরির ধাপ কীভাবে চলবে
}

\vspace{0.5mm}

\textbf{
Programming Model
=
Program-এর Code/Structure কীভাবে সাজানো হবে
}

\end{center}

\end{itemize}

\end{rememberbox}

\small

%================================================
% 5.17 Modular Programming
%================================================
